\documentclass[a4paper,11pt]{ltjsarticle}

\usepackage{amsmath}
\usepackage{amssymb}
\usepackage{bm}
\usepackage{mathtools}
\usepackage{graphicx}
\usepackage{hyperref}
\usepackage{physics}
\usepackage{mathrsfs}


\begin{document}
\section{一様重力場中の質点の運動}
鉛直上向きを $z$ 軸の正方向とする。質量 $m$ の質点が一様な重力場中を運動している。
重力加速度を $g$ とすると、運動方程式は
\begin{equation}
m\ddot{z}=-mg
\label{eq.motion}
\end{equation}
である。したがって、
\[
\ddot{z}=-g
\]
となる。初期条件を
\[
z(0)=z_{0} \quad \dot{z}(0)=v_{0}
\]
とすると、速度と位置は
\begin{align*}
\dot{z}(t)&=v_{0}-gt\\
z(t)&=z_{0}+v_{0}t-\frac{1}{2}gt^2
\end{align*}
となる。また、ポテンシャルエネルギーを $V(z)=mgz$ とすれば
、この系のラグランジアンは
\begin{equation}
\mathcal{L}=\frac{1}{2}m\dot{z}^2-mgz
\end{equation}

\begin{equation}
\frac{d}{dt}\pdv{\mathcal{L}}{\dot{z}}-\pdv{\mathcal{L}}{z}=0
\label{eq.Euler-Lagrange}
\end{equation}
を用いると、同じ運動方程式\eqref{eq.motion}が得られる。

\subsection{エネルギー保存}
ラグランジアンからではなく、まずNewton方程式\eqref{eq.motion}
の両辺に $\dot{z}$ を掛け、
\[
m\dot{z}\ddot{z}=-mg\dot{z}
\]
から
\begin{equation}
\frac{d}{dt}\left(\frac{1}{2}m\dot{z}^2+mgz\right)=0
\end{equation}
が得られる。

ラグランジアンは $\mathcal{L}$ で表し、ラグランジアン密度は $\mathscr{L}$ で表す。


\end{document}